\begin{thebibliography}{43}
\bibitem{ref1}Noga Alon, Jeff Edmonds, and Michael Luby. Linear time erasure codes with nearly optimal recovery. In Proceedings of IEEE 36th Annual Foundations of Computer Science, pages 512–519. IEEE, 1995.
\bibitem{ref2}Noga Alon and Michael Luby. A linear time erasure-resilient code with nearly optimal recovery. IEEE Transactions on Information Theory, 42(6):1732–1736, 1996.
\bibitem{ref3}Sanjeev Arora, Carsten Lund, Rajeev Motwani, Madhu Sudan, and Mario Szegedy. Proof verification and the hardness of approximation problems. Journal of the ACM (JACM), 45(3):501–555, 1998.
\bibitem{ref4}Sanjeev Arora and Shmuel Safra. Probabilistic checking of proofs: A new characterization of NP. Journal of the ACM (JACM), 45(1):70–122, 1998.
\bibitem{ref5}László Babai, Lance Fortnow, Leonid A Levin, and Mario Szegedy. Checking computations in polylogarithmic time. In Proceedings of the twenty-third annual ACM symposium on Theory of computing, pages 21–32, 1991.
\bibitem{ref6}Laszlo Babai, Lance Fortnow, Noam Nisan, and Avi Wigderson. BPP has subexponential time simulations unless EXPTIME has publishable proofs. Computational Complexity, 3(4):307–318, 1993.
\bibitem{ref7}Omer Barkol, Yuval Ishai, and Ronny Roth. Locally decodable codes and their applications. PhD thesis, Computer Science Department, Technion, 2008.
\bibitem{ref8}Donald Beaver and Joan Feigenbaum. Hiding instances in multioracle queries. In Annual Symposium on Theoretical Aspects of Computer Science, pages 37–48. Springer, 1990.
\bibitem{ref9}Mihir Bellare and John Rompel. Randomness-efficient oblivious sampling. In Proceedings of the 35th Annual Symposium on Foundations of Computer Science, 1994, pages 276–287. IEEE, 1994.
\bibitem{ref10}Manuel Blum, Michael Luby, and Ronitt Rubinfeld. Self-testing/correcting with applications to numerical problems. In Proceedings of the twenty-second annual ACM symposium on Theory of computing, pages 73–83, 1990.
\bibitem{ref11}Benny Chor, Oded Goldreich, Eyal Kushilevitz, and Madhu Sudan. Private information retrieval. In Proceedings of IEEE 36th Annual Foundations of Computer Science, pages 41–50. IEEE, 1995.
\bibitem{ref12}Zeev Dvir. On matrix rigidity and locally self-correctable codes. computational complexity, 20(2):367–388, 2011.
\bibitem{ref13}Zeev Dvir, Parikshit Gopalan, and Sergey Yekhanin. Matching vector codes. SIAM Journal on Computing, 40(4):1154–1178, 2011.
\bibitem{ref14}Klim Efremenko. 3-query locally decodable codes of subexponential length. SIAM Journal on Computing, 41(6):1694–1703, 2012.
\bibitem{ref15}Peter Gemmell, Richard Lipton, Ronitt Rubinfeld, Madhu Sudan, and Avi Wigderson. Self- testing/correcting for polynomials and for approximate functions. In STOC, volume 91, pages 32–42. Citeseer, 1991.
\bibitem{ref16}Peter Gemmell and Madhu Sudan. Highly resilient correctors for polynomials. Information processing letters, 43(4):169–174, 1992.
\bibitem{ref17}Oded Goldreich. A sample of samplers: A computational perspective on sampling. In Studies in Complexity and Cryptography, pages 302–332. Springer, 2011. doi:10.1007/ 978-3-642-22670-0\_24.
\bibitem{ref18}Oded Goldreich and Madhu Sudan. Locally testable codes and PCPs of almost-linear length. Journal of the ACM (JACM), 53(4):558–655, 2006.
\bibitem{ref19}Sivakanth Gopi, Swastik Kopparty, Rafael Oliveira, Noga Ron-Zewi, and Shubhangi Saraf. Locally testable and locally correctable codes approaching the gilbert-varshamov bound. IEEE Transactions on Information Theory, 64(8):5813–5831, 2018.
\bibitem{ref20}Alan Guo, Swastik Kopparty, and Madhu Sudan. New affine-invariant codes from lifting. In Proceedings of the 4th conference on Innovations in Theoretical Computer Science, pages 529–540. ACM, 2013.
\bibitem{ref21}Venkatesan Guruswami, Atri Rudra, and Madhu Sudan. Essential coding theory. Draft available at http://www.cse.buffalo.edu/\textasciitilde{}atri/courses/coding-theory/book, 2012.
\bibitem{ref22}Brett Hemenway, Rafail Ostrovsky, and Mary Wootters. Local correctability of expander codes. Information and Computation, 243:178–190, 2015.
\bibitem{ref23}Jonathan Katz and Luca Trevisan. On the efficiency of local decoding procedures for error- correcting codes. In Proceedings of the thirty-second annual ACM symposium on Theory of computing, pages 80–86, 2000.
\bibitem{ref24}Tali Kaufman and Madhu Sudan. Sparse random linear codes are locally decodable and testable. In 48th Annual IEEE Symposium on Foundations of Computer Science (FOCS’07), pages 590–600. IEEE, 2007.
\bibitem{ref25}Kiran S Kedlaya and Sergey Yekhanin. Locally decodable codes from nice subsets of finite fields and prime factors of mersenne numbers. SIAM Journal on Computing, 38(5):1952–1969, 2009.
\bibitem{ref26}Swastik Kopparty, Or Meir, Noga Ron-Zewi, and Shubhangi Saraf. High-rate locally correctable and locally testable codes with sub-polynomial query complexity. Journal of the ACM (JACM), 64(2):11, 2017.
\bibitem{ref27}Swastik Kopparty, Shubhangi Saraf, and Sergey Yekhanin. High-rate codes with sublinear-time decoding. Journal of the ACM (JACM), 61(5):28, 2014.
\bibitem{ref28}Ray Li and Mary Wootters. Lifted multiplicity codes and the disjoint repair group property. In Approximation, Randomization, and Combinatorial Optimization. Algorithms and Techniques (APPROX/RANDOM 2019). Schloss Dagstuhl-Leibniz-Zentrum fuer Informatik, 2019.
\bibitem{ref29}Richard J. Lipton. Efficient checking of computations. In Annual Symposium on Theoretical Aspects of Computer Science, pages 207–215. Springer, 1990.
\bibitem{ref30}Irving S Reed. A class of multiple-error-correcting codes and the decoding scheme. Technical report, Massachusetts inst of tech Lexington Lincoln lab, 1953.
\bibitem{ref31}Omer Reingold, Salil Vadhan, and Avi Wigderson. Entropy waves, the zig-zag graph product, and new constant-degree expanders and extractors. In Electronic Colloquium on Computational Complexity (ECCC), page 18, 2001. URL: https://eccc.weizmann.ac.il/report/2001/018/.
\bibitem{ref32}Ronitt Rubinfeld and Madhu Sudan. Robust characterizations of polynomials with applications to program testing. SIAM Journal on Computing, 25(2):252–271, 1996.
\bibitem{ref33}C. E. Shannon. A mathematical theory of communication. ACM SIGMOBILE Mobile Computing and Communications Review, 5(1):3–55, 2001. Originally appeared in Bell System Tech. J. 27:379–423, 623–656, 1948.
\bibitem{ref34}Carl Siegel. Über die classenzahl quadratischer zahlkörper. Acta Arithmetica, 1(1):83–86, 1935.
\bibitem{ref35}Madhu Sudan, Luca Trevisan, and Salil Vadhan. Pseudorandom generators without the XOR lemma. Journal of Computer and System Sciences, 62(2):236–266, 2001.
\bibitem{ref36}Luca Trevisan. List-decoding using the XOR lemma. In 44th Annual IEEE Symposium on Foundations of Computer Science, 2003. Proceedings., pages 126–135. IEEE, 2003.
\bibitem{ref37}Arnold Walfisz. Zur additiven zahlentheorie. ii. Mathematische Zeitschrift, 40(1):592–607, 1936.
\bibitem{ref38}David Woodruff. New lower bounds for general locally decodable codes. In Electronic Colloquium on Computational Complexity (ECCC), volume 14, 2007.
\bibitem{ref39}S. Yekhanin. Locally decodable codes. In International Computer Science Symposium in Russia, pages 289–290. Springer, 2011.
\bibitem{ref40}Sergey Yekhanin. Towards 3-query locally decodable codes of subexponential length. Journal of the ACM (JACM), 55(1):1–16, 2008.
\bibitem{ref41}Kalina Petrova Zeev Dvir. Lecture 1: Introduction. Lecture notes: https://www.cs. princeton.edu/\textasciitilde{}zdvir/LDCnotes/LDC1.pdf, year = 2016.
\bibitem{ref42}Kalina Petrova Zeev Dvir. Lecture 4: Lower bounds for r-query LDCs. Lecture notes: https://www.cs.princeton.edu/\textasciitilde{}zdvir/LDCnotes/LDC4.pdf, 2016.
\bibitem{ref43}Victor Vasilievich Zyablov. An estimate of the complexity of constructing binary linear cascade codes. Problemy Peredachi Informatsii, 7(1):5–13, 1971.
\end{thebibliography}
